\documentclass[12pt]{article}

\usepackage{amsmath}
\usepackage{amssymb}
\usepackage{amsfonts}
\usepackage{geometry}
\usepackage{setspace}
\usepackage[T1]{fontenc}
\usepackage[utf8]{inputenc}


\geometry{margin=0.8in}
\setstretch{1.2}

\begin{document}

\subsection*{Kou (2002): A Jump-Diffusion Model for Option Pricing}

\subsubsection*{Introduction}

The Black--Scholes model assumes log-normally distributed asset returns, but two
well-documented empirical failures motivate Kou (2002). First, real return distributions
exhibit a \textit{leptokurtic} shape: a higher peak and heavier, asymmetric tails than
the normal distribution. Second, implied volatilities backed out from market prices are
not flat across strikes but form a convex \textit{volatility smile}, meaning Black--Scholes
systematically underprices options with extreme strikes. Kou proposes a double exponential
jump-diffusion model that addresses both problems while retaining analytical tractability
for a wide class of derivatives.

\subsubsection*{The Double Exponential Jump-Diffusion Model}

The stock price is governed by:

\begin{equation}
    \frac{dS(t)}{S(t^{-})} = \mu\,dt + \sigma\,dW(t)
    + d\!\left(\sum_{i=1}^{N(t)}(V_i - 1)\right)
\end{equation}

where $W(t)$ is a standard Brownian motion, $N(t)$ a Poisson process with intensity
$\lambda$, and $\{V_i\}$ i.i.d.\ non-negative jump multipliers independent of $W(t)$ and
$N(t)$. The log-jump size $Y = \log V$ follows an asymmetric double exponential
distribution:

\begin{equation}
    f_Y(y) = p\,\eta_1 e^{-\eta_1 y}\mathbf{1}_{\{y\geq 0\}}
           + q\,\eta_2 e^{\eta_2 y}\mathbf{1}_{\{y < 0\}},
    \qquad \eta_1>1,\;\eta_2>0
\end{equation}

where $p+q=1$ are the probabilities of upward and downward jumps, with mean jump sizes
$1/\eta_1$ and $1/\eta_2$ respectively. The constraint $\eta_1>1$ ensures finite expected
stock price. By choosing $q>p$ and calibrating $\eta_2$ independently of $\eta_1$, the
model captures the empirically observed asymmetry between large negative and positive
moves --- something the symmetric normal jumps in Merton (1976) cannot reproduce.

\subsubsection*{Why Double Exponential: Intuition and Tractability}

The double exponential distribution is motivated on two grounds. Economically, it links
to behavioral finance: the heavy tails model investor \textit{overreaction} to news
(large jumps), while the high peak models \textit{underreaction} (small price responses).
The asymmetry between upward and downward components directly explains the volatility
skew, since higher probability of large downward jumps inflates implied volatility for
low-strike puts.

Mathematically, the key advantage is the \textit{memoryless property} of the exponential
distribution. When a jump process crosses a barrier, it often overshoots. Computing the
exact distribution of this overshoot is generally intractable, but for exponential jump
sizes the memoryless property gives a closed-form overshoot distribution and ensures the
overshoot is conditionally independent of the first passage time. This is precisely why
the Kou model yields analytical solutions for barrier, lookback, and perpetual American
options, while the Merton model with normal jumps cannot.

For European options, the pricing formula uses the $Hh$ function:

\begin{equation}
    Hh_n(x) = \frac{1}{n!}\int_x^{\infty}(t-x)^n e^{-t^2/2}\,dt,
    \qquad n = 0,1,2,\ldots
\end{equation}

which satisfies a simple three-term recursion and is computed efficiently from the
standard normal distribution. In practice, truncating the series to ten to fifteen
terms gives accurate results.

\subsubsection*{Strengths, Limitations, and Importance}

The model's principal strengths are its analytical tractability across vanilla and
path-dependent derivatives, its asymmetric jump structure that fits the volatility skew,
and the clear economic interpretation of all parameters ($\lambda$ is mean jumps per year;
$p, q$ are jump direction probabilities; $1/\eta_i$ are mean jump sizes). It improves on
Merton (1976) both empirically and mathematically, and it can be embedded in a rational
expectations equilibrium, making its pricing formula preference-free.

The main limitations mirror those of the Merton model. The constant jump intensity
$\lambda$ cannot capture time-varying crash risk, and the Poisson independence assumption
rules out volatility clustering. The constant diffusion volatility $\sigma$ ignores the
well-documented mean-reversion of return volatility; models combining stochastic
volatility with jumps, such as Bates (1996), address this. The market remains incomplete
under jump dynamics, so additional equilibrium assumptions are needed to uniquely pin
down prices.

\subsubsection*{Relevance to QQQ Options}

The QQQ ETF tracks the NASDAQ-100, a technology-heavy index prone to large, sudden
price moves driven by earnings surprises, regulatory decisions, and product announcements.
The QQQ implied volatility surface shows a pronounced downside skew, with
out-of-the-money puts carrying materially higher implied volatilities than at-the-money
options. The Kou model's asymmetric jump structure is well suited to match this pattern
by assigning higher probability and larger size to downward jumps through the parameters
$q$ and $\eta_2$. We calibrate $(\sigma, \lambda, p, \eta_1, \eta_2)$ to the QQQ option
cross-section and compare the fit against Black--Scholes and Merton. The Kou model is
expected to outperform both benchmarks for short-maturity, far-from-the-money options, especially on the put side, while converging toward the Merton fit for longer-dated near-the-money contracts where the diffusion component dominates.

\end{document}
